\part{Expériences et études de cas}

Cette partie ne contient aucune performance inventée. Chaque chapitre commence comme protocole. Les tableaux de résultats ne sont remplis qu'après exécution avec provenance et statut.

\chapter{EXP-001 — Benchmark de coordination}
\section{All-to-all, registry et capability routing}
Comparer les mêmes traces de tâches et de changements sous cinq architectures : couplage pairwise explicite, graphe statique, registry simple, graphe centralisé et dynamic capability graph.

\section{Complexité, latence et maintenance}
Mesurer relations déclarées, modifications de configuration, ruptures de contrats, latence de résolution, coût de calcul et interventions humaines. La relation $N(N-1)/2$ est utilisée uniquement comme baseline structurelle du graphe complet.

\section{Analyse de sensibilité}
Faire varier $N$, le nombre de capacités, la multiplicité des fournisseurs, le churn de versions et la proportion de tâches inter-repos.

\chapter{EXP-002 — Substitution de fournisseurs}
\section{Implémentations concurrentes}
Choisir plusieurs capacités disposant d'au moins deux providers ou construire des providers contrôlés équivalents pour la mécanique initiale.

\section{Contrats Capability IR}
Comparer types sémantiques, unités, qualité, evidence, latence et coût avant/après substitution.

\section{Robustesse aux changements}
Mesurer modifications des consommateurs, adaptations, violations de contrats et temps de récupération. Les substitutions incompatibles servent de contrôles négatifs.

\chapter{EXP-003 — Divergence théorie-code}
\section{Injection de divergences}
Créer des variantes commit-addressées contenant des défauts d'API, d'unité, convention, baseline, contrôle négatif, incertitude, provenance, hypothèse et conclusion.

\section{Tests classiques, CI sémantique, CI scientifique et OAK}
Geler les variantes puis exécuter séparément les détecteurs. Les runners ne reçoivent pas le label du défaut avant que leurs receipts soient figés.

\section{Taux de détection et faux positifs}
Comparer détection par famille, contrôles propres, temps de détection, coût de calcul et effort de revue. Un gain de détection ne peut être interprété sans le coût correspondant.

\chapter{EXP-004 — Reuse-first contre generate-first}
\section{Protocoles expérimentaux}
Comparer une stratégie generate-first à `SEARCH → REUSE → ADAPT → GENERATE` sur un ensemble de tâches gelé et randomisé ou apparié.

\section{Duplication et temps de développement}
Mesurer implémentations sémantiquement redondantes, nouveaux modules, symboles réutilisés, lignes modifiées, temps et appels d'outils.

\section{Qualité et maintenabilité}
Mesurer tests, défauts, revue humaine et coût des changements ultérieurs. L'historique seul n'est pas utilisé pour attribuer un effet causal.

\chapter{EXP-005 — Étude empirique du GitHub Tristan}
\section{Corpus et provenance}
Le corpus initial inclut `TFUGA-AI7-TRISTAN2` et `TTM-TFUGA-AI7-TRISTAN2`, avec expansion seulement vers des dépôts autorisés et réellement pertinents.

\section{Chronologie des capacités}
Construire les lignées Repository→PR→Commit→File/Symbol→Capability→Claim→Evidence.

\section{Réutilisation, M-moins et cristallisation}
Mesurer apparition, réutilisation, duplication, convergence, échecs conservés, promotion et dérive théorie-code. Les résultats restent observationnels.

\chapter{EXP-006 — Mycelial Fixed-Point Reconstruction}
\section{Seed minimal}
Geler une version de l'écosystème et définir exactement quels schémas, manifests, graphes, tests et compilateurs sont conservés.

\section{Reconstruction clean-room}
Reconstruire les artefacts dérivés dans un espace isolé sans dépendances implicites autorisées hors du seed déclaré.

\section{Score de fermeture}
Mesurer fractions de capacités, tests, claims, evidence et programmes récupérés, ainsi que les dépendances cachées. Le verdict porte sur la reproductibilité architecturale, pas sur la vérité scientifique.

\chapter{EXP-007 — Theory-Code-Evidence Round Trip}
\section{Théorie de départ}
Choisir une théorie bornée avec définitions, hypothèses, claims, unités/types et limites explicites.

\section{Compilation et evidence}
Compiler les obligations, produire l'implémentation ou les HOLD explicites, puis reconstruire la théorie implicite à partir des artefacts observables.

\section{Résidu de round-trip}
Comparer définitions, hypothèses, claims, unités/types, evidence, limites et provenance. Un faible résidu mesure la cohérence interne dans le protocole; il ne prouve pas le contenu scientifique.

\chapter{Échelle de preuve expérimentale}
\section{Niveaux}
Les campagnes progressent selon
\[
Protocol\to SyntheticMechanics\to HistoricalReplay\to ProspectiveObservation\to CleanRoomReplay\to IndependentReplication.
\]
Aucun niveau n'est déduit automatiquement du précédent.

\section{Rapports négatifs}
Tout FAIL ou HOLD conserve ses paramètres, traces et conséquences dans M$^-$ ou M$^?$. Les expériences négatives ne sont pas supprimées pour améliorer la conclusion.
